\documentclass[10pt]{article}
\usepackage[notes]{aditya}
\usepackage{extarrows}

\title{MAT 545 Notes}
\author{Aditya Pahuja}
\date{\today}

\begin{document}
\maketitle
\tableofcontents
\pagebreak

\section{08/24}

\emph{In which we vaguely discuss the concepts that will be presented
during the semester.}


{\scriptsize Maybe I will revise this section once the semester ends.}

\subsection{Textbook}

We will cover the first two chapters (0 and 1) of
\emph{Principles of Algebraic Geometry} by Griffiths and Harris.
Aleksey Zinger has a
\href{https://www.math.stonybrook.edu/~azinger/mat545-fall19/GHnotes.pdf}{list of typos/corrections}
for the book (link downloads a pdf).

\subsection{Course overview}

In what follows, green bold words are vocabulary words
that I for the most part don't know/care to define right now.

The goal of this course is to study complex manifolds,
and in particular \vocab{projective} and \vocab{K\"ahler manifolds}.
Ultimately, we will build up to a proof of the
\alert{Kodaira embedding theorem}, which says that
if the \vocab{K\"ahler form} $\omega$ of a compact K\"ahler manifold $X$
is \vocab{rational}, then $X$ is projective, meaning
it can be embedded in some (large enough) projective space.
% One reason we care about projective manifolds is that
% projective spaces have many sub-projective spaces,
% whose intersections with an embedded $X$, if transverse,
% give many natural subvarieties/submanifolds to get our hands on
% when trying to understand $X$'s geometry.

\subsubsection{Several complex variables}
If we want to work with complex manifolds, we will have to do
analysis in several complex variables, analogous to the real case.
Many things from single variable complex analysis carry over
to the $n$-variable case; here is a small list.
\begin{itemize}
    \ii Taylor series.
    \ii The Cauchy integral formula.
    \ii Identity theorem: if a holomorphic function
    is zero on an open subset of its domain, then it is identically zero
    (but NOT the stronger version, in which
    a holomorphic function on $U\subseteq\CC$ is zero identically
    if its zero set has an accumulation point).
    \ii Some extension phenomena if the holomorphic function is bounded\todo{?}
    \ii Weierstrass preparation theorem:
    if $f\colon U\to\CC$ is holomorphic on an open
    $U\subseteq\CC\times\CC^{n-1}$,
    then $f(w,z) = w^k + f_{k-1}(z)w^{k-1} + \cdots + f_0(z)$
    where the $f_i$ are holomorphic on $\CC^{n-1}$,
    in some small neighborhood of zero.
    When $n=1$ this manifests as $f(w) = w^k$ locally.
    \todo{say more about WPT (zero set branched cover blah blah\dots?)}
\end{itemize}

On the other hand, there are phenomena in higher dimensions
that are entirely different from the single variable case.
For example, \vocab{Hartog's theorem} says that
for $n\ge 2$, any holomorphic function on the punctured $n$-disk in $\CC^n$
admits a holomorphic extension to the whole disk;
this is obviously not true when $n = 1$, since $f(z) = \frac 1z$
is unbounded near zero.

\subsubsection{Tangent, cotangent bundles, cohomology}

Let $M$ be a complex manifold of (complex) dimension $n$.
The real tangent bundle $T^\RR M$ is a rank $2n$ real vector bundle,
and its complexification is the rank $2n$ complex vector bundle
$T^\CC M \coloneq T^\RR M\otimes_\RR \CC = T^\RR M\oplus i T^\RR M$.
This complexification splits canonically
into a \vocab{holomorphic part} $T^{1,0}M$
and a \vocab{antihomolomorphic part} $T^{0,1}M$,
both rank $n$ complex vector bundles;
the holomorphic part is spanned by the vectors
$\frac\partial{\partial z_j}$
and the antiholomorphic part is spanned by the vectors
$\frac\partial{\partial\ol z_j}$.
The holomorphic and antiholomorphic cotangent bundles
$\check T^{1,0}M$ and $\check T^{0,1}M$
are then defined as the duals of the respective tangent bundles
and are respectively spanned by covectors of the form
$\dd z_j$ and $\dd\ol z_j$.

We can also define complex differential $k$-forms as
elements of $\Lambda^k (\check T^\RR\otimes\CC)$,
which are written in local coordinates as
holomorphic linear combinations of wedges of $\dd z_i$ and $\dd\ol z_i$.
It turns out that the number of covectors in the holomorphic cotangent bundle
and the number of covectors in the antiholomorphic tangent bundle
needed to express an addend of a $k$-form in local coordinates
is independent of the coordinates chosen, so we get a splitting
\begin{equation*}
    \Lambda^k(\check T^\RR M\otimes\CC) = \bigoplus_{p+q=k}\check T^{p,q}M
\end{equation*}
for each $k$, and these $\check T^{p,q}M$ are functorially defined.
One can then define complex de~Rham cohomology in the usual way,
and \vocab{Dolbeault cohomology} $H^{p,q}_{\ol\partial}$,
which looks at $\ol\partial$-closed and exact forms
in the space of $(p,q)$-forms.
These two cohomology theories are related by the
\vocab{Hodge-de~Rham spectral sequence}
(which we will not cover in this course).

We will also study \vocab{sheaves},
which package ``local'' information in a way that
lets us try to extract ``global'' information (often manifesting as
sheaf cohomology). For example, we will prove
the de~Rham and Dolbeault theorems via the machinery of sheaves.

\subsubsection{Connections and Chern classes}
For a holomorphic vector bundle $E$, we can define
analogues of standard differential-geometry structures, such as
connections, curvature, and metrics.
To this end, let $h$ be a \vocab{Hermitian metric} on $E$.
We will show that there is a canonical \vocab{connection} $D_h$
on $E$ associated with $h$, analogous to the Levi-Civita connection.
Any connection then gives rise to a notion of \vocab{curvature} $D^2$.

In the case of line bundles, $D^2$ turns out to be a
$2$-form with coefficients in the bundle $\hom(E,E)$;
using the tensor-hom adjunction we can show that this bundle is trivial,
so that we can treat it as a bona fide $2$-form.
In fact, it is a closed $(1,1)$-form, and the associated cohomology class,
up to scaling by $2\pi i$, belongs to the integral cohomology $H^2(M;\ZZ)$.
Moreover, this cohomology class is the \vocab{first Chern class} of $E$!

\subsubsection{Hodge theory}

Let $M$ be a compact K\"ahler manifold with K\"ahler form $\omega$.
Let $\mathcal H^k_\omega(M)$ be the space of $k$-forms on $M$
which are \vocab{harmonic} with respect to $\omega$.
It turns out that
\begin{equation*}
    H^k(M;\RR)\cong\mathcal H^k_\omega(M)\cong\mathcal H^{p,q}_\omega(M)
    \cong H^{p,q}(M),
\end{equation*}
which is the \vocab{Hodge decomposition}.\todo{is the coef field right?}
This is one tool that can be used to prove the following:
\begin{theorem}[Hard Lefschetz theorem]
    \label{thm:hard-lefschetz}
    Let $M$ be an $n$-dimensional compact K\"ahler manifold
    with K\"ahler form $\omega$. Then, there is an isomorphism
    \begin{equation*}
	H^{n-k}(M;\RR)\overset{\omega^k\smile}\too H^{n+k}(M;\RR).
    \end{equation*}
\end{theorem}

\begin{theorem}[Weak Lefschetz theorem]
    \label{thm:weak-lefschetz}
    If $N^{n-1}\into M^n$ is a hypersurface, then
    \todo{probably need more conditions on $N$}
    \begin{equation*}
        i^*\colon H^j(M;\ZZ)\to H^j(N;\ZZ)
    \end{equation*}
    is an isomorphism for $j<n-1$ and injective for $j=n-1$.
\end{theorem}
\todo{idk i don't have much else to say}

\end{document}
